% =============================================================
% CAPÍTULO 4 — Resultados de compresión SVD
% Alimentado por: NB02 (espectral, uniforme, adaptativa) y
%                 NB03 (sensibilidad por componente y profundidad)
% =============================================================

\section{Resultados de compresión SVD}
\label{sec:resultados-compresion}

Este capítulo analiza el comportamiento de BERT-base \emph{fine-tuned} sobre GoEmotions
(23 emociones) bajo distintas estrategias de compresión basadas en la descomposición
en valores singulares (SVD). El objetivo es doble: primero, determinar cuánta
información se puede eliminar sin destruir la capacidad emocional del modelo, y segundo,
identificar dónde reside esa información dentro de la arquitectura.

El capítulo se organiza en cuatro secciones. La primera examina la estructura espectral
de las 72 matrices lineales del modelo (\S\ref{sec:analisis-espectral}). La segunda
evalúa la compresión uniforme y documenta una transición de fase abrupta
(\S\ref{sec:compresion-uniforme}). La tercera desglosa la sensibilidad por tipo de
componente y por profundidad (\S\ref{sec:sensibilidad-componente}). La cuarta presenta
la compresión adaptativa como alternativa informada por la estructura espectral
(\S\ref{sec:compresion-adaptativa}).

El modelo de referencia (\emph{baseline}) a lo largo de todo el capítulo es BERT-base
\emph{fine-tuned} sin ninguna compresión, con un F1 macro de 0{,}577 sobre el conjunto de test.


% =============================================================
% §4.1 — Análisis espectral
% =============================================================

\subsection{Análisis espectral de las matrices de peso}
\label{sec:analisis-espectral}

Antes de aplicar cualquier compresión, conviene entender la estructura interna de los
pesos del modelo. Cada capa de BERT-base contiene seis matrices lineales: Query ($W_Q$),
Key ($W_K$), Value ($W_V$) y Attention Output ($W_O$) en el bloque de self-attention, y
FFN Intermediate ($W_{\text{inter}}$) y FFN Output ($W_{\text{out}}$) en el bloque
feed-forward. Con 12 capas, el modelo contiene 72 matrices lineales en total.

Para cada matriz $W \in \mathbb{R}^{m \times n}$, calculamos su descomposición en valores
singulares $W = U \Sigma V^\top$ y definimos la energía acumulada al rango $k$
como la fracción de la norma de Frobenius capturada por los primeros $k$ valores singulares:

\begin{equation}
    E(k) = \frac{\sum_{i=1}^{k} \sigma_i^2}{\sum_{i=1}^{n} \sigma_i^2}
    \label{eq:energia-acumulada}
\end{equation}

El rango mínimo $k_\tau$ para un umbral de energía $\tau$ es el menor $k$ tal
que $E(k) \geq \tau$. Un valor bajo de $k_\tau$ indica que la matriz concentra su
información en pocos componentes principales y es, por tanto, más compresible.

\subsubsection{Asimetría self-attention vs.\ feed-forward}

La Figura~\ref{fig:anatomy} muestra el mapa de rangos mínimos $k_{95}$ para cada una de
las 72 matrices, organizadas por capa (eje vertical) y tipo de componente (eje horizontal).
El resultado más inmediato es una asimetría pronunciada entre los componentes de
self-attention y los de feed-forward.

\begin{figure}[H]
    \centering
    \includegraphics[width=0.85\textwidth]{figures/cap4/cap4_anatomy_heatmap_es.png}
    \caption{Rango mínimo $k_{95}$ por capa y tipo de componente. Colores fríos indican alta compresibilidad; colores cálidos, baja. Se observa una división clara entre los componentes de self-attention (izquierda) y los feed-forward (derecha).}
    \label{fig:anatomy}
\end{figure}

La Tabla~\ref{tab:rangos-medios} cuantifica esta diferencia. Las matrices Query y Key
requieren un rango medio de aproximadamente 395 para retener el 95\% de la energía
espectral, mientras que las matrices FFN Intermediate y Output necesitan rangos superiores
a 618. El gap medio entre Q y FFN Intermediate es de unos 222 rangos, lo que equivale
al 29\% del rango completo (768).

\begin{table}[H]
    \centering
    \caption{Rango medio $k_{95}$ por tipo de componente (media sobre las 12 capas).}
    \label{tab:rangos-medios}
    \small
    \begin{tabular}{@{}lcc@{}}
        \toprule
        \textbf{Componente} & $\boldsymbol{\mu(k_{95})}$ & \textbf{Compresibilidad} \\
        \midrule
        Query (Q)           & $\sim$396 & Alta \\
        Key (K)             & $\sim$394 & Alta \\
        Value (V)           & $\sim$424 & Alta \\
        Attn Output         & $\sim$434 & Media-alta \\
        FFN Intermediate    & $\sim$618 & Baja \\
        FFN Output          & $\sim$625 & Baja \\
        \bottomrule
    \end{tabular}
\end{table}

Esta asimetría es consistente a lo largo de las 12 capas: no se trata de un artefacto
puntual de una o dos capas anómalas, sino de una propiedad estructural del modelo
\emph{fine-tuned}. La interpretación es directa: las matrices Q y K concentran los patrones de
atención en un subespacio de dimensionalidad efectiva mucho menor que 768, mientras que
las matrices FFN distribuyen su información de forma más uniforme a lo largo de todo el
espectro de valores singulares.

\subsubsection{Perfil de decaimiento espectral}

La Figura~\ref{fig:fingerprint} refuerza esta observación desde una perspectiva
complementaria. Muestra las curvas de decaimiento de los valores singulares (media $\pm$
desviación estándar sobre las 12 capas) para cada tipo de componente.

\begin{figure}[H]
    \centering
    \includegraphics[width=0.85\textwidth]{figures/cap4/cap4_spectral_fingerprint_es.png}
    \caption{Curvas de decaimiento espectral por tipo de componente (media $\pm$ std
    sobre 12 capas). Q y K presentan decaimiento rápido; FFN Intermediate y Output
    mantienen valores singulares elevados a lo largo de todo el espectro.}
    \label{fig:fingerprint}
\end{figure}

Las matrices Q y K exhiben un decaimiento pronunciado: los primeros componentes principales
capturan la mayor parte de la varianza y los restantes contribuyen marginalmente. Las
matrices FFN, en cambio, presentan un espectro esencialmente plano, donde cada valor
singular adicional aporta una fracción similar de información. Esta diferencia cualitativa
en la forma del espectro explica por qué un mismo rango de truncamiento tiene efectos
radicalmente distintos sobre ambos tipos de componente.

\subsubsection{Solapamiento de subespacios Q-K-V}

Como observación complementaria, analizamos el solapamiento entre los subespacios
dominantes de las matrices Q, K y V dentro de cada capa, medido como la similitud de
coseno media entre los vectores singulares principales.

\begin{table}[H]
    \centering
    \caption{Rango de solapamiento entre subespacios (top-32 y top-64 vectores singulares).}
    \label{tab:overlap}
    \small
    \begin{tabular}{@{}lcc@{}}
        \toprule
        \textbf{Par} & \textbf{Top-32} & \textbf{Top-64} \\
        \midrule
        Q--K & 0{,}34 -- 0{,}63 & 0{,}31 -- 0{,}47 \\
        Q--V & 0{,}03 -- 0{,}18 & 0{,}06 -- 0{,}25 \\
        K--V & 0{,}05 -- 0{,}21 & 0{,}08 -- 0{,}27 \\
        \bottomrule
    \end{tabular}
\end{table}

Q y K comparten un subespacio significativo (solapamiento de hasta 0{,}63 para los 32
vectores singulares principales), lo cual es esperable dado que su producto $Q \cdot K^\top$
genera los pesos de atención. V, en cambio, opera en un subespacio sustancialmente
ortogonal a los de Q y K. Este resultado es consistente con la diferencia funcional de
ambos bloques: Q y K definen \emph{a qué atender}, mientras que V define \emph{qué
información extraer}. La Figura~\ref{fig:qkv-overlap} visualiza esta asimetría
capa a capa.

\begin{figure}[H]
    \centering
    \includegraphics[width=0.85\textwidth]{figures/cap4/cap4_qkv_overlap_es.png}
    \caption{Solapamiento de subespacios entre Q, K y V por capa, medido como la
    similitud de coseno media entre los \emph{top}-32 vectores singulares de cada par.
    El par Q--K (azul) presenta solapamiento sistemáticamente alto en todas las
    capas, coherente con que $QK^\top$ genere los pesos de atención. Los pares
    Q--V y K--V (naranja, verde) operan en subespacios prácticamente ortogonales.}
    \label{fig:qkv-overlap}
\end{figure}

No explotamos este solapamiento para compresión conjunta, pero lo
documentamos como hallazgo descriptivo de la geometría interna del modelo.

\subsubsection{Estructura espectral de las activaciones internas}
 
El análisis anterior se ha centrado en la geometría de las matrices de
peso. Una pregunta complementaria es cómo evoluciona la estructura
espectral de las \emph{activaciones} a lo largo de la red, es decir, qué
tan concentrada está la información que sale efectivamente de cada capa
durante la inferencia. La Figura~\ref{fig:cap4-internal-compression}
documenta tres métricas calculadas sobre las activaciones del modelo
\emph{fine-tuneado}, evaluadas sobre el conjunto de test.
 
\begin{figure}[H]
    \centering
    \includegraphics[width=0.95\textwidth]{figures/cap4/cap4_internal_compression_es.png}
    \caption{Caracterización espectral de las activaciones internas por
    capa. \textbf{Arriba, eje izquierdo (azul):} norma euclidiana media
    de los vectores de salida, separada en tokens de contenido y token
    \texttt{[CLS]}. \textbf{Arriba, eje derecho (rojo):} rango efectivo
    derivado de la entropía espectral de la matriz de activaciones, y
    umbral $k_{95}$ (número mínimo de valores singulares que captura el
    95\% de la energía). El sombreado de fondo distingue tres regímenes:
    léxico (Emb--L3), mezcla (L4--L7) y colapso (L8--L11).
    \textbf{Abajo:} heatmap de energía espectral acumulada por capa
    (filas) y posición del valor singular ordenado de mayor a menor
    (columnas). La línea de puntos blancos marca $k_{95}$. En L4--L7
    hacen falta $\sim 250$ valores singulares para capturar el 95\% de
    la energía; en L11, menos de 50.}
    \label{fig:cap4-internal-compression}
\end{figure}
 
La estructura es informativa para el problema de compresión. Las capas
intermedias (L4--L7) operan en una geometría de rango efectivo alto
(en torno a 200), lo que indica que la información se distribuye sobre
muchas direcciones independientes; la redundancia es alta y el espacio
está bien aprovechado. Las capas tardías (L8--L11), en cambio, presentan
un colapso espectral progresivo: el rango efectivo cae a aproximadamente
30 en L11, y el $k_{95}$ se reduce desde $\sim 250$ en L7 hasta menos de
50 al final del modelo. Esto significa que las capas tardías han
\emph{comprimido} ya su información en muy pocas direcciones.
 
Esta caracterización proporciona una explicación estructural del gradiente
de fragilidad que se documentará en \S\ref{sec:sensibilidad-componente}:
truncar a rango 128 una capa con rango efectivo 200 elimina información
que el modelo estaba usando, pero conserva la mayor parte de la señal;
truncar a rango 128 una capa con rango efectivo 30 es matemáticamente
inocuo (el modelo no usa más de 30 direcciones de todas formas), pero
la señal residual es tan concentrada que cualquier perturbación destruye
predicciones específicas. La capacidad de las capas tardías está
\emph{ya comprimida internamente}, y eso las hace estructuralmente
imposibles de comprimir más sin romperlas.


\subsubsection{Predicciones a partir del análisis espectral}

Los resultados espectrales generan tres predicciones concretas que las secciones
siguientes contrastarán experimentalmente:

\begin{enumerate}[nosep]
    \item Las matrices Q y K deberían tolerar una reducción de rango agresiva (hasta
    $r \approx 128$) con degradación mínima del rendimiento, dado que su dimensionalidad
    efectiva es muy inferior a 768.

    \item Las matrices FFN deberían ser el cuello de botella de cualquier estrategia de
    compresión, al requerir rangos sustancialmente mayores para retener una fracción
    equivalente de energía.

    \item Una estrategia de compresión uniforme, que aplique el mismo rango a todos los
    componentes, será necesariamente subóptima, al ignorar esta heterogeneidad espectral.
\end{enumerate}

Las secciones \S\ref{sec:compresion-uniforme} y \S\ref{sec:sensibilidad-componente}
confirman las dos primeras predicciones, y la \S\ref{sec:compresion-adaptativa} valida
la tercera.


% =============================================================
% §4.2 — Compresión uniforme
% =============================================================

\subsection{Compresión uniforme y transición de fase}
\label{sec:compresion-uniforme}

La estrategia más simple de compresión SVD consiste en aplicar el mismo rango de
truncamiento $r$ a todas las matrices lineales del modelo simultáneamente. Esta sección
evalúa cinco niveles de rango uniforme ($r \in \{512, 384, 256, 128, 64\}$) y documenta
una transición de fase abrupta en el rendimiento.

\subsubsection{Degradación global}

La Tabla~\ref{tab:uniform} muestra el F1 macro, el ratio de parámetros y
la retención de rendimiento para cada nivel de rango.

\begin{table}[H]
    \centering
    \caption{Resultados de compresión uniforme sobre el conjunto de test.
    El \emph{ratio de parámetros} se define como
    $\text{params}_{\text{comprimido}}/\text{params}_{\text{original}}$:
    valores $<1$ indican compresión real; $=1$, sin reducción; $>1$,
    expansión paramétrica.}
    \label{tab:uniform}
    \small
    \begin{tabular}{@{}lccc@{}}
        \toprule
        \textbf{Estrategia} & \textbf{F1 macro} & \textbf{Ratio params.} & \textbf{Retención} \\
        \midrule
        Baseline           & 0{,}577 & 1{,}00 & 100{,}0\% \\
        Uniform $r$=512    & 0{,}464 & 1{,}00 & 80{,}5\%  \\
        Uniform $r$=384    & 0{,}251 & 0{,}81 & 43{,}5\%  \\
        Uniform $r$=256    & 0{,}025 & 0{,}61 & 4{,}3\%   \\
        Uniform $r$=128    & 0{,}000 & 0{,}42 & 0{,}0\%   \\
        Uniform $r$=64     & 0{,}000 & 0{,}32 & 0{,}0\%   \\
        \bottomrule
    \end{tabular}
\end{table}

El patrón dominante es una \emph{transición de fase} entre $r=384$ y $r=256$: el modelo
pasa de retener el 43\% del rendimiento a retener apenas el 4\%, una caída de casi 40
puntos porcentuales en un solo salto de rango. A partir de $r=128$, el modelo entra en
``muerte clínica'' con F1 exactamente igual a 0{,}0: no es capaz de clasificar
correctamente ni una sola emoción. La Figura~\ref{fig:uniform-transition} ilustra
visualmente este \emph{cliff}.

\begin{figure}[H]
    \centering
    \includegraphics[width=0.78\textwidth]{figures/cap4/cap4_uniform_transition_es.png}
    \caption{Transición de fase bajo compresión uniforme. F1 macro
    en función del rango: la curva se mantiene en niveles funcionales
    para $r \geq 384$ y colapsa abruptamente entre $r=384$ y $r=256$. La
    región sombreada marca la zona de muerte clínica ($F1 = 0$).}
    \label{fig:uniform-transition}
\end{figure}

La zona útil de compresión uniforme se limita a $r=512$ (degradación significativa pero
funcional) y $r=384$ (degradación severa pero no total). Cualquier rango inferior destruye
la capacidad del modelo de forma irrecuperable.

\subsubsection{Vulnerabilidad emocional diferencial}

La degradación no afecta a todas las emociones por igual. La Figura~\ref{fig:emotion-vulnerability}
muestra el F1 por emoción para cada nivel de compresión uniforme, revelando un espectro
amplio de resistencia.

\begin{figure}[H]
    \centering
    \includegraphics[width=0.75\textwidth]{figures/cap4/cap4_emotion_vulnerability_es.png}
    \caption{F1 por emoción para cada nivel de compresión uniforme. Las emociones están
    ordenadas por su F1 en el baseline. Las emociones frecuentes y léxicamente simples
    (arriba) sobreviven más que las infrecuentes y semánticamente complejas (abajo).}
    \label{fig:emotion-vulnerability}
\end{figure}

Las emociones más resistentes comparten características comunes: son frecuentes en el
corpus y están asociadas a marcadores léxicos fuertes. \emph{Gratitude}, por ejemplo,
retiene un F1 de 0{,}92 a $r=384$ y aún mantiene 0{,}55 a $r=256$, apoyada por la presencia
recurrente de expresiones como ``thank you'' y ``thanks''. \emph{Amusement} (0{,}74 a
$r=384$) y \emph{love} (0{,}77 a $r=384$) siguen un patrón similar.

En el extremo opuesto, \emph{embarrassment} colapsa ya a $r=512$ (de 0{,}27 a 0{,}00),
\emph{realization} cae de 0{,}27 a 0{,}16 en el mismo nivel, y \emph{disappointment} pasa
de 0{,}28 a 0{,}06. Estas emociones son infrecuentes en el corpus de entrenamiento,
semánticamente ambiguas y carecen de marcadores léxicos distintivos. Su representación
depende de dimensiones espectrales de cola, precisamente las primeras en ser eliminadas
por la compresión SVD.

Este patrón conecta directamente con los resultados de cristalización del
Capítulo~\ref{sec:resultados-interpretabilidad}: las emociones que cristalizan temprano en
la profundidad del modelo (es decir, que son reconocibles desde las primeras capas) tienden
a ser las más resistentes a la compresión, mientras que las que requieren procesamiento
profundo son las más vulnerables. La \S\ref{sec:sintesis} formaliza esta correlación.


% =============================================================
% §4.3 — Sensibilidad por componente y profundidad
% =============================================================

\subsection{Sensibilidad por componente y profundidad}
\label{sec:sensibilidad-componente}

La sección anterior mostró que la compresión uniforme presenta una transición de fase
abrupta. El análisis espectral de la \S\ref{sec:analisis-espectral} predijo que los
distintos tipos de componente deberían tener sensibilidades radicalmente diferentes.
Esta sección lo verifica experimentalmente mediante un estudio sistemático de 33
configuraciones, aislando el efecto del tipo de componente (8 tipos $\times$ 3 rangos)
y de la profundidad (3 grupos $\times$ 3 rangos).

\subsubsection{Sensibilidad por tipo de componente}

Para aislar el efecto de cada tipo de componente, comprimimos selectivamente solo las 12
matrices de un tipo dado (por ejemplo, las 12 matrices Query, una por capa) a un rango
fijo, manteniendo el resto del modelo intacto. La Tabla~\ref{tab:component-sensitivity}
presenta los resultados al rango $r=128$, que maximiza el contraste entre componentes.

\begin{table}[H]
    \centering
    \caption{Sensibilidad por componente a $r=128$. La sensibilidad normalizada mide
    los puntos porcentuales de F1 perdidos por cada 1\% de parámetros eliminados.}
    \label{tab:component-sensitivity}
    \small
    \begin{tabular}{@{}lcccc@{}}
        \toprule
        \textbf{Componente} & \textbf{F1 macro} & \textbf{Retención} & \textbf{Sensibilidad} \\
        \midrule
        Query (Q)           & 0{,}574 & 99{,}4\% & 0{,}14 \\
        Key (K)             & 0{,}568 & 98{,}4\% & 0{,}38 \\
        Value (V)           & 0{,}395 & 68{,}5\% & 7{,}31 \\
        Attn Output         & 0{,}325 & 56{,}3\% & 10{,}13 \\
        FFN Output          & 0{,}128 & 22{,}2\% & 3{,}80 \\
        FFN Intermediate    & 0{,}040 & 6{,}9\%  & 4{,}55 \\
        Attention (todas)   & 0{,}027 & 4{,}6\%  & --- \\
        FFN (todas)         & 0{,}001 & 0{,}1\%  & --- \\
        \bottomrule
    \end{tabular}
\end{table}

El resultado confirma y amplifica la predicción espectral. A rango 128, Query retiene el
99,4\% del rendimiento mientras que FFN Intermediate retiene apenas el 6,9\%: un factor
de 14$\times$ en retención absoluta. Normalizado por la fracción de parámetros
eliminados, la diferencia en sensibilidad es de 72$\times$ (0,14 frente a 10,13
para Attn Output, el componente de atención más sensible).

Este resultado valida la primera predicción del análisis espectral: Q y K son
efectivamente ``inmunes'' a la compresión agresiva, porque su información reside en un
subespacio de dimensionalidad mucho menor que el rango completo de la matriz.

\subsubsection{Tres regímenes de compresibilidad}

Al examinar la evolución de la retención a lo largo de los tres rangos evaluados
($r \in \{256, 128, 64\}$), los seis tipos de componente se agrupan en tres regímenes
cualitativamente distintos:

\begin{description}[nosep, leftmargin=1.2em]
    \item[Régimen~1 --- Inmunes (Q, K):] Degradación gradual y casi lineal. Query
    pierde apenas 1{,}8 puntos porcentuales entre $r=256$ y $r=64$. Estas matrices
    operan en un subespacio de dimensionalidad efectiva muy inferior a 768.

    \item[Régimen~2 --- Efecto cliff (V, Attn Output):] Sobreviven a $r=256$
    (retención $\sim$93\%) pero colapsan a $r=64$ ($\sim$3--5\%). Existe un umbral
    crítico entre $r=128$ y $r=64$ por debajo del cual la información no se puede
    recuperar. Value, por ejemplo, cae 63 puntos porcentuales solo entre $r=128$ y
    $r=64$.

    \item[Régimen~3 --- Frágiles (FFN Inter, FFN Output):] Ya gravemente dañados a
    $r=256$ (52\% y 75\% de retención respectivamente). El cliff principal se sitúa
    entre $r=256$ y $r=128$. Son matrices de alta dimensionalidad efectiva cuya
    información está distribuida a lo largo de todo el espectro de valores singulares.
\end{description}

\begin{figure}[H]
    \centering
    \includegraphics[width=0.85\textwidth]{figures/cap4/cap4_component_sensitivity_es.png}
    \caption{Curvas de retención de F1 en función del rango para cada uno
    de los seis tipos de componente. Los tres regímenes
    descritos en el texto se separan visualmente: Q y K (azules) decaen
    suavemente; V y Attn Output (naranjas) presentan el \emph{cliff}
    entre $r=128$ y $r=64$; las dos FFN (rojas) caen abruptamente desde
    $r=256$. La línea horizontal punteada marca el 50\% de retención.}
    \label{fig:component-sensitivity}
\end{figure}

Estos tres regímenes corresponden directamente a las tres formas de espectro observadas
en la \S\ref{sec:analisis-espectral}: decaimiento rápido (Q, K), decaimiento intermedio
(V, Attn Output) y espectro plano (FFN). La estructura espectral predice con precisión
el comportamiento bajo compresión.

\subsubsection{Sensibilidad por profundidad}

Para evaluar el efecto de la profundidad, comprimimos simultáneamente todas las matrices
de un grupo de capas (Early: 0--3, Middle: 4--7, Late: 8--11) manteniendo el resto del
modelo intacto.

\begin{table}[H]
    \centering
    \caption{Retención de F1 por grupo de profundidad y rango de compresión.}
    \label{tab:depth-sensitivity}
    \small
    \begin{tabular}{@{}lccc@{}}
        \toprule
        \textbf{Grupo} & $\boldsymbol{r=256}$ & $\boldsymbol{r=128}$ & $\boldsymbol{r=64}$ \\
        \midrule
        Early (0--3)   & 85{,}4\% & 62{,}5\% & 37{,}0\% \\
        Middle (4--7)  & 82{,}5\% & 46{,}5\% & 25{,}3\% \\
        Late (8--11)   & 33{,}3\% &  0{,}0\% &  0{,}0\% \\
        \bottomrule
    \end{tabular}
\end{table}

Las capas tardías (8--11) son \emph{catastróficamente sensibles}: a $r=128$, el modelo
muere por completo (F1 = 0{,}0). En contraste, las capas tempranas retienen más del 60\%
del rendimiento al mismo rango. La Figura~\ref{fig:depth-sensitivity} sintetiza este
gradiente.

\begin{figure}[H]
    \centering
    \includegraphics[width=0.85\textwidth]{figures/cap4/cap4_depth_sensitivity_es.png}
    \caption{Sensibilidad por grupo de profundidad. Cada barra representa
    la retención de F1 macro al comprimir todas las matrices de un grupo
    de capas a un rango dado, manteniendo el resto intacto. Las capas
    tardías (rojo) son las únicas que colapsan por completo a $r=128$;
    las tempranas (verde) y medias (azul) retienen capacidad funcional.}
    \label{fig:depth-sensitivity}
\end{figure}

Este resultado es consistente con la literatura de
interpretabilidad mecánica, que identifica las capas finales de los Transformers como las
responsables de la representación semántica de alto nivel. En el contexto de la clasificación
emocional, las capas 8--11 son las que codifican las distinciones finas entre emociones;
destruir su estructura equivale a destruir la capacidad clasificadora.

Este hallazgo complementa el análisis por componente: la compresión uniforme falla no solo
porque trata igual a Q y a FFN, sino porque trata igual a las capas 0 y a las capas 11.
Una estrategia de compresión verdaderamente informada debería respetar ambos ejes de
asimetría: tipo de componente y profundidad.

\subsubsection{Nota metodológica: silhouette scores en contextos de colapso}
\label{sec:silhouette-caveat}

Como métrica complementaria al F1, evaluamos los silhouette scores sobre
proyecciones t-SNE de los embeddings del modelo comprimido. Los resultados
revelan una limitación importante de esta métrica en contextos de
compresión severa que conviene documentar como aporte metodológico
transferible a otros estudios de compresión.

Los silhouette scores negativos observados para FFN Intermediate a $r=64$
($-0{,}125$) y $r=128$ ($-0{,}051$) indican que la compresión no solo
destruye la estructura de clusters sino que la \emph{desordena}
activamente: las muestras quedan más cerca de clusters ajenos que del
propio.

Sin embargo, las capas tardías a $r=128$ presentan un silhouette score de
$0{,}15$, el más alto de todas las configuraciones, a pesar de un F1
de $0{,}0$. Esto constituye un \emph{artefacto sistemático}: cuando el
modelo colapsa, todos los embeddings convergen a un único cluster trivial,
y la métrica de silhouette, que premia la cohesión intra-cluster y la
separación inter-cluster, interpreta erróneamente esa convergencia
patológica como ``buena estructura''. Documentamos esta observación como
una limitación general de las métricas geométricas de clustering en
contextos de colapso del modelo, y recomendamos que cualquier estudio de
compresión que utilice silhouette o métricas relacionadas (Davies-Bouldin,
Calinski-Harabasz) verifique siempre la validez de la métrica mediante
inspección directa de las proyecciones cuando el F1 sea bajo. La regla
operativa es simple: si el modelo colapsa por F1, las métricas de
clustering geométrico no son interpretables.


% =============================================================
% §4.4 — Compresión adaptativa
% =============================================================

\subsection{Compresión adaptativa y recetas mixtas}
\label{sec:compresion-adaptativa}

Los resultados de las secciones anteriores establecen que la compresión uniforme es
inherentemente subóptima: aplica el mismo rango a componentes con sensibilidades que
difieren en un factor de hasta 72$\times$ (normalizado por parámetros eliminados) y
a capas con vulnerabilidades que van del 0\% al 62\% de retención al mismo rango. La compresión adaptativa propone una alternativa:
asignar a cada matriz un rango proporcional a su complejidad espectral, determinado por
el umbral de energía $\tau$.

\subsubsection{Resultados de compresión adaptativa}

En la compresión adaptativa, cada matriz recibe el rango mínimo $k_\tau$ necesario para
retener una fracción $\tau$ de su energía espectral. Matrices con espectro concentrado
(Q, K) reciben rangos bajos; matrices con espectro plano (FFN) reciben rangos altos. La
Tabla~\ref{tab:adaptive} presenta los resultados para cinco umbrales de energía.

\begin{table}[H]
    \centering
    \caption{Resultados de compresión adaptativa y recetas mixtas.
    \emph{Ratio params.} =
    $\text{params}_{\text{comprimido}}/\text{params}_{\text{original}}$,
    donde $<1$ es compresión real y $>1$ expansión paramétrica
    (ver \S\ref{sec:caveat-compresion}).}
    \label{tab:adaptive}
    \small
    \begin{tabular}{@{}lccc@{}}
        \toprule
        \textbf{Estrategia} & \textbf{F1 macro} & \textbf{Ratio params.} & \textbf{Retención} \\
        \midrule
        Baseline             & 0{,}577 & 1{,}00 & 100{,}0\% \\
        Adaptive 99\%        & 0{,}558 & 1{,}20 & 96{,}7\%  \\
        Mixed Conservative   & 0{,}505 & 1{,}04 & 87{,}5\%  \\
        Adaptive 95\%        & 0{,}491 & 1{,}02 & 85{,}1\%  \\
        Mixed Aggressive     & 0{,}382 & 0{,}91 & 66{,}3\%  \\
        Adaptive 90\%        & 0{,}357 & 0{,}89 & 61{,}8\%  \\
        Adaptive 85\%        & 0{,}260 & 0{,}80 & 45{,}1\%  \\
        Adaptive 80\%        & 0{,}143 & 0{,}72 & 24{,}7\%  \\
        \bottomrule
    \end{tabular}
\end{table}

La compresión adaptativa al 99\% de energía retiene el 96.7\% del rendimiento, con
una degradación de apenas 1{,}9 puntos porcentuales de F1 respecto al baseline. Al 95\%,
la retención es del 85.1\%, lo cual constituye un punto de equilibrio razonable.

Las recetas mixtas combinan compresión adaptativa con reglas específicas por tipo de
componente (por ejemplo, comprimir más agresivamente Q y K mientras se protegen las FFN).
La receta conservadora (87.5\% de retención) se sitúa entre Adaptive 95\% y 99\%,
ofreciendo un compromiso intermedio.

\subsubsection{Frontera de Pareto: adaptativa vs.\ uniforme}

La Figura~\ref{fig:pareto} compara las dos familias de estrategias en el espacio
retención--compresión. La compresión adaptativa \emph{domina} a la uniforme en toda la
frontera de Pareto: para cualquier nivel de retención deseado, la adaptativa lo consigue
con un ratio de compresión igual o superior.

\begin{figure}[H]
    \centering
    \includegraphics[width=0.85\textwidth]{figures/cap4/cap4_pareto_frontier_es.png}
    \caption{Frontera de Pareto: compresión uniforme (cuadrados) vs.\ adaptativa
    (círculos). La adaptativa domina en todo el rango. Los puntos Mixed se sitúan
    sobre la frontera adaptativa.}
    \label{fig:pareto}
\end{figure}

La ventaja de la compresión adaptativa es más pronunciada en la zona de alta retención
(>80\%), donde la uniforme solo dispone de un punto útil ($r=512$) mientras que la
adaptativa ofrece un continuo de opciones entre 85\% y 97\% de retención.

\subsubsection{Distribución de rangos asignados}

La Figura~\ref{fig:rank-distribution} ilustra cómo la compresión adaptativa distribuye
los rangos entre los distintos tipos de componente para cada umbral de energía.

\begin{figure}[H]
    \centering
    \includegraphics[width=0.85\textwidth]{figures/cap4/cap4_adaptive_rank_distribution_es.png}
    \caption{Distribución de rangos asignados por la compresión adaptativa para cada
    umbral de energía, desglosada por tipo de componente. Q y K reciben rangos
    consistentemente menores que FFN.}
    \label{fig:rank-distribution}
\end{figure}

La distribución confirma la lógica espectral: a un umbral del 95\%, Q y K reciben rangos
del orden de 390--400, mientras que FFN Intermediate y Output reciben rangos de 610--625.
La adaptativa simplemente \emph{respeta} la estructura natural de las matrices en lugar
de forzar un rango uniforme.

\subsubsection{Caveat: ratio de compresión y simplificación de rango}
\label{sec:caveat-compresion}

Es necesario documentar una limitación importante de la compresión adaptativa a umbrales
altos. La factorización SVD truncada a rango $r$ reemplaza una matriz $W \in \mathbb{R}^{m
\times n}$ por dos matrices $U_r \in \mathbb{R}^{m \times r}$ y $V_r \in \mathbb{R}^{r
\times n}$ (absorbiendo $\Sigma_r$ en una de las dos). El número de parámetros pasa de
$mn$ a $r(m + n)$. Para que exista compresión real, se requiere $r < \frac{mn}{m+n}$.

En el caso de BERT-base, las matrices de self-attention tienen dimensiones $768 \times 768$,
lo que da un umbral de $r < 384$ para compresión real. Las matrices FFN tienen dimensiones
$768 \times 3072$ (Intermediate) y $3072 \times 768$ (Output), con un umbral de
$r < \frac{768 \times 3072}{768 + 3072} = 614.4$.

Las estrategias Adaptive 95\% y 99\% asignan, por tanto, rangos que para muchas
matrices \emph{superan} el umbral de compresión real. Los ratios de parámetros de
1{,}02 (Adaptive 95\%) y 1{,}20 (Adaptive 99\%) confirman que estas estrategias
producen un modelo con \emph{más} parámetros que el original.

En rigor, estas estrategias no constituyen compresión de parámetros sino
\emph{simplificación de rango}: la representación factorizada tiene menor rango
matricial que la original, lo cual puede ser útil para análisis e interpretabilidad, pero
no reduce el tamaño del modelo. La compresión real solo comienza a partir de Adaptive
90\% (ratio 0{,}89), y los niveles con compresión significativa (85\% y 80\%)
ya implican degradaciones del 55\% y 75\% respectivamente.


% =============================================================
% Cierre del capítulo
% =============================================================

\subsection{Síntesis del capítulo}
\label{sec:sintesis-compresion}

Los resultados de este capítulo establecen cuatro hallazgos principales:

\begin{enumerate}[nosep]
    \item \textbf{Asimetría espectral.} Las matrices de self-attention (Q, K) concentran
    su información en un subespacio de dimensionalidad efectiva $\sim$400, mientras que
    las matrices FFN requieren $\sim$620 dimensiones. Este gap de $\sim$220 rangos es
    una propiedad estructural del modelo \emph{fine-tuned}, consistente en las 12 capas.

    \item \textbf{Transición de fase en compresión uniforme.} La compresión uniforme
    presenta un cliff abrupto entre $r=384$ (43\% de retención) y $r=256$ (4\%).
    La zona útil se limita a $r \geq 384$, y cualquier rango inferior destruye el
    modelo de forma irrecuperable.

    \item \textbf{Sensibilidad heterogénea.} Los componentes presentan sensibilidades
    que difieren en un factor de 72$\times$ (normalizado por parámetros). Q y K son
    inmunes a compresión agresiva; V y Attn Output presentan un cliff entre $r=128$ y
    $r=64$; las FFN son frágiles desde $r=256$. Las capas tardías (8--11) son
    catastróficamente sensibles, con muerte total a $r=128$.

    \item \textbf{Dominancia matizada de la compresión adaptativa.} La
    compresión adaptativa supera a la uniforme \emph{en el rango donde hay
    compresión real} (umbrales del 80\% y 90\%, con ratios de parámetros
    de 0{,}72 y 0{,}89 respectivamente). Por encima del 90\%, los umbrales
    del 95\% y 99\% producen ratios de parámetros $>1$ ($1{,}02$ y
    $1{,}20$ respectivamente): no son técnicamente compresión, sino
    \emph{simplificación de rango con expansión paramétrica}. La aparente
    dominancia de la adaptativa en toda la frontera de Pareto debe
    interpretarse, por tanto, con esta cautela: cuando el ratio de
    parámetros es genuinamente menor que el original ($<1$), la adaptativa
    supera a la uniforme; en la región $\geq 1$, ambas familias se
    encuentran fuera del régimen de compresión efectiva.
\end{enumerate}

Estos hallazgos motivan dos líneas de investigación que se desarrollan en los capítulos
siguientes: la interpretabilidad mecánica (Capítulo~\ref{sec:resultados-interpretabilidad}),
que busca explicar \emph{por qué} ciertos componentes y capas son más críticos, y la
síntesis (\S\ref{sec:sintesis}), que conecta la vulnerabilidad a compresión con la
profundidad de cristalización emocional.